% Sección incluida desde ARTURITO MAIN.tex con % Sección incluida desde ARTURITO MAIN.tex con % Sección incluida desde ARTURITO MAIN.tex con % Sección incluida desde ARTURITO MAIN.tex con \input{explicando_firmware.tex}.
% Este archivo contiene un fragmento del informe, no un documento independiente.

% =============================================================================
\section{Diseño del firmware, simulación y depuración}
\label{sec:diseno-firmware}
% =============================================================================

\subsection{Arquitectura propuesta y criterios de diseño}

El firmware se diseñará para integrar la percepción, la estimación de la pose,
la construcción del mapa y el control del rover diferencial. La arquitectura
será portable entre plataformas: el desarrollo y la validación funcional se
realizarán primero sobre ESP32; luego se portará a EDU-CIAA reutilizando la
aplicación y los drivers compartidos, y reemplazando la capa de adaptación al
hardware.

Se adopta una organización en tres capas para separar las operaciones propias
de cada placa de los algoritmos del robot. Esta separación permite probar los
módulos por partes, mantener una única implementación de la lógica del rover y
evitar que las decisiones de una placa condicionen el código común. La
arquitectura de ejecución se organizará en tareas periódicas controladas por
tiempo y planificadas por FreeRTOS. El scheduler será apropiativo y asignará el
procesador según las prioridades; las tareas se bloquearán durante sus
intervalos de espera. Por lo tanto, el modelo es de tareas periódicas con
planificación apropiativa, no un superloop cooperativo. Se utilizará FreeRTOS
en ambas plataformas.

\begin{center}
\fbox{\parbox{0.96\linewidth}{\centering
Sensores y actuadores $\leftrightarrow$ Capa 1: HAL de la placa
$\leftrightarrow$ Capa 2: drivers de sensores y motores
$\leftrightarrow$ Capa 3: adquisición, pose, mapa y navegación
$\leftrightarrow$ comunicaciones $\leftrightarrow$ interfaz de PC}}
\end{center}

El modelo elegido prioriza interfaces estables entre capas, unidades físicas
comunes y un runtime compartido. El período de muestreo y las prioridades se
ajustarán mediante mediciones de carga y latencia cuando se integren los
encoders, la construcción del mapa y el control de motores.

\subsection{Pseudocódigo del funcionamiento previsto}

\begin{enumerate}
  \item Configurar la placa, el scheduler, el canal de telemetría y el estado
        inicial seguro de los actuadores.
  \item Inicializar la IMU MPU6050, el sensor ToF VL53L0X, los encoders AS5600
        y el puente H DRV8833. Informar qué periféricos no están disponibles;
        no habilitar el movimiento si falta una condición necesaria para el
        control seguro.
  \item Calibrar los sensores que lo requieran y poner en cero la pose inicial
        $(x,y,\theta)$. Preparar la representación del mapa correspondiente:
        acumulación polar para V0 o grilla de ocupación para V1.
  \item Iniciar las tareas periódicas de FreeRTOS para adquisición,
        procesamiento de la aplicación y publicación/recepción de datos.
  \item En cada ciclo de adquisición, leer gyro Z, distancia ToF y ángulos de
        ambas ruedas; asociar las muestras con el instante de lectura y
        entregarlas a la aplicación.
  \item Procesar las muestras: corregir el gyro, estimar el desplazamiento de
        cada rueda, actualizar la pose del rover y validar la distancia.
  \item Actualizar el mapa con el rayo medido por el ToF y la pose vigente.
        En V0 se acumularán mediciones en un plot polar suponiendo que el rover
        gira en el origen; en V1 se actualizará una grilla de ocupación 2D en
        coordenadas mundo usando $(x,y,\theta)$.
  \item Según el modo seleccionado, obtener la consigna desde navegación
        autónoma o desde teleoperación, convertirla en velocidades de rueda y
        comandar el DRV8833 mediante la HAL. Ante una condición de parada o una
        pérdida de comunicación que supere el límite definido, detener los
        motores.
  \item Publicar pose, estado y datos del mapa para la interfaz externa;
        repetir el ciclo mientras el sistema esté activo.
\end{enumerate}

La navegación autónoma, el control de motores, la integración de AS5600 y la
grilla V1 forman parte del desarrollo previsto; no se presentan como
funcionalidades ya verificadas. SLAM permanece como objetivo secundario: la
primera grilla se construirá con la pose estimada por odometría.

\subsection{Módulos organizados en capas}

\begin{description}
  \item[Capa 1 --- HAL de periféricos.] Tendrá una interfaz común y una
        implementación por placa. Inicializará los periféricos y entregará
        magnitudes en unidades acordadas: velocidad angular en rad/s,
        distancia en metros, ángulos de rueda y señales de mando para motores.
        Incluirá HAL para tiempo, MPU6050, VL53L0X, AS5600, GPIO/PWM del DRV8833
        y telemetría. No realizará integración de pose, mapeo ni navegación.
  \item[Capa 2 --- drivers compartidos.] Interpretará las señales físicas de
        cada dispositivo y entregará datos listos para la aplicación. Incluirá
        la calibración y corrección del giroscopio, validación de distancia,
        lectura y conversión de ángulos AS5600, y generación de señales de
        control para el DRV8833. Sus interfaces no dependerán de Arduino ni de
        SAPI.
  \item[Capa 3 --- aplicación.] Coordinará las tareas FreeRTOS y contendrá la
        lógica propia del rover: adquisición y sincronización de muestras,
        cinemática diferencial y estimación de pose, actualización de la grilla
        de ocupación, selección de consignas según el modo y distribución de
        datos a la navegación y las comunicaciones. Los algoritmos de mapa y
        navegación no se ubicarán en la HAL ni en los drivers de dispositivos.
\end{description}

El runtime común organizará el trabajo en tareas de adquisición, procesamiento
de drivers/aplicación y publicación de telemetría. La división interna de la
tarea de procesamiento podrá refinarse al medir los tiempos de ejecución; el
contrato entre capas se mantendrá igual para las dos placas.

\subsection{Bibliotecas y herramientas externas previstas}

\begin{itemize}
  \item \textbf{ESP32:} framework Arduino y su integración con FreeRTOS;
        \texttt{Wire} para I2C; biblioteca Adafruit \texttt{VL53L0X} para el
        sensor ToF. Para el transporte inalámbrico se utilizará la capacidad
        disponible en ESP32 una vez elegido el protocolo de enlace.
  \item \textbf{EDU-CIAA:} LPCOpen y SAPI para la placa y sus periféricos, más
        el kernel FreeRTOS configurado en el proyecto. La HAL de ToF implementa
        el acceso I2C por registros y no utiliza la biblioteca Adafruit.
  \item \textbf{Código del proyecto:} contratos HAL, drivers de MPU6050,
        VL53L0X, AS5600 y DRV8833, y módulos de pose, mapa y navegación. Los
        drivers de AS5600 y DRV8833 se incorporarán al contar con el hardware y
        sus interfaces acordadas.
  \item \textbf{PC:} Processing y su biblioteca Serial se utilizan en V0 para
        visualizar las muestras recibidas. La interfaz objetivo mostrará la
        pose del rover y el mapa actualizado.
\end{itemize}

\subsection{Software de PC e interfaz de usuario prevista}

El alcance confirmado para el programa de PC es visualizar la pose del rover y
el mapa actualizado. El enlace serie USB/UART se utiliza en V0; para transmitir
pose y mapa sin cable se prevé incorporar Wi-Fi o Bluetooth en una etapa
posterior. La elección entre ambos protocolos queda pendiente. Los
requerimientos generales también mencionan órdenes START/STOP y, como objetivos
secundarios, teleoperación y planificación de rutas; su incorporación como
controles de esta interfaz queda por definir.

\begin{center}
\fbox{\parbox{0.91\linewidth}{\centering
Inicio $\rightarrow$ conexión con el rover $\rightarrow$ espera de datos
$\rightarrow$ visualización continua de pose y mapa $\rightarrow$ pérdida de
enlace: indicar desconexión / conexión restablecida: reanudar visualización}}
\end{center}

Si se pierde el enlace durante la operación, el firmware deberá aplicar la
parada segura y la interfaz indicará la desconexión. El flujo representa la
interfaz prevista; las pantallas de Processing disponibles actualmente solo
reciben y dibujan telemetría.

\subsection{Ensayos realizados y estado de validación}

Los ensayos se realizaron con el firmware V0 y hardware conectado, usando
programas de prueba y Processing; no se dispone de una simulación integral del
rover. Se separaron perfiles de prueba para compilar la lectura de IMU y la del
ToF de manera independiente. También se implementó \texttt{imu\_compare}, que
transmite gyro Z crudo, ángulo crudo integrado y ángulo procesado para comparar
la integración sin corrección con la que aplica calibración y zona muerta.

El firmware V0 de ESP32 fue cargado y validado con Processing: se comprobó la
recepción de las muestras serie y la visualización de orientación y distancia.
El contrato de aplicación V0 transmite a 115200 baudios dos columnas,
\texttt{theta\_rad,distancia\_m}; una distancia inválida se representa con
$-1$. El build actual con FreeRTOS compartido compila, pero su telemetría debe
confirmarse después de la última modificación. El firmware EDU-CIAA compila y
enlaza, aunque falta probarlo con sensores conectados a la placa.

Los archivos disponibles no contienen capturas de pantalla ni resultados
numéricos de precisión o error. Para documentar visualmente los ensayos se debe
incorporar una captura de Processing durante la prueba V0 y, cuando se realice,
otra del gráfico de comparación de la IMU. Estas evidencias no se reemplazan
por imágenes ilustrativas.
.
% Este archivo contiene un fragmento del informe, no un documento independiente.

% =============================================================================
\section{Diseño del firmware, simulación y depuración}
\label{sec:diseno-firmware}
% =============================================================================

\subsection{Arquitectura propuesta y criterios de diseño}

El firmware se diseñará para integrar la percepción, la estimación de la pose,
la construcción del mapa y el control del rover diferencial. La arquitectura
será portable entre plataformas: el desarrollo y la validación funcional se
realizarán primero sobre ESP32; luego se portará a EDU-CIAA reutilizando la
aplicación y los drivers compartidos, y reemplazando la capa de adaptación al
hardware.

Se adopta una organización en tres capas para separar las operaciones propias
de cada placa de los algoritmos del robot. Esta separación permite probar los
módulos por partes, mantener una única implementación de la lógica del rover y
evitar que las decisiones de una placa condicionen el código común. La
arquitectura de ejecución se organizará en tareas periódicas controladas por
tiempo y planificadas por FreeRTOS. El scheduler será apropiativo y asignará el
procesador según las prioridades; las tareas se bloquearán durante sus
intervalos de espera. Por lo tanto, el modelo es de tareas periódicas con
planificación apropiativa, no un superloop cooperativo. Se utilizará FreeRTOS
en ambas plataformas.

\begin{center}
\fbox{\parbox{0.96\linewidth}{\centering
Sensores y actuadores $\leftrightarrow$ Capa 1: HAL de la placa
$\leftrightarrow$ Capa 2: drivers de sensores y motores
$\leftrightarrow$ Capa 3: adquisición, pose, mapa y navegación
$\leftrightarrow$ comunicaciones $\leftrightarrow$ interfaz de PC}}
\end{center}

El modelo elegido prioriza interfaces estables entre capas, unidades físicas
comunes y un runtime compartido. El período de muestreo y las prioridades se
ajustarán mediante mediciones de carga y latencia cuando se integren los
encoders, la construcción del mapa y el control de motores.

\subsection{Pseudocódigo del funcionamiento previsto}

\begin{enumerate}
  \item Configurar la placa, el scheduler, el canal de telemetría y el estado
        inicial seguro de los actuadores.
  \item Inicializar la IMU MPU6050, el sensor ToF VL53L0X, los encoders AS5600
        y el puente H DRV8833. Informar qué periféricos no están disponibles;
        no habilitar el movimiento si falta una condición necesaria para el
        control seguro.
  \item Calibrar los sensores que lo requieran y poner en cero la pose inicial
        $(x,y,\theta)$. Preparar la representación del mapa correspondiente:
        acumulación polar para V0 o grilla de ocupación para V1.
  \item Iniciar las tareas periódicas de FreeRTOS para adquisición,
        procesamiento de la aplicación y publicación/recepción de datos.
  \item En cada ciclo de adquisición, leer gyro Z, distancia ToF y ángulos de
        ambas ruedas; asociar las muestras con el instante de lectura y
        entregarlas a la aplicación.
  \item Procesar las muestras: corregir el gyro, estimar el desplazamiento de
        cada rueda, actualizar la pose del rover y validar la distancia.
  \item Actualizar el mapa con el rayo medido por el ToF y la pose vigente.
        En V0 se acumularán mediciones en un plot polar suponiendo que el rover
        gira en el origen; en V1 se actualizará una grilla de ocupación 2D en
        coordenadas mundo usando $(x,y,\theta)$.
  \item Según el modo seleccionado, obtener la consigna desde navegación
        autónoma o desde teleoperación, convertirla en velocidades de rueda y
        comandar el DRV8833 mediante la HAL. Ante una condición de parada o una
        pérdida de comunicación que supere el límite definido, detener los
        motores.
  \item Publicar pose, estado y datos del mapa para la interfaz externa;
        repetir el ciclo mientras el sistema esté activo.
\end{enumerate}

La navegación autónoma, el control de motores, la integración de AS5600 y la
grilla V1 forman parte del desarrollo previsto; no se presentan como
funcionalidades ya verificadas. SLAM permanece como objetivo secundario: la
primera grilla se construirá con la pose estimada por odometría.

\subsection{Módulos organizados en capas}

\begin{description}
  \item[Capa 1 --- HAL de periféricos.] Tendrá una interfaz común y una
        implementación por placa. Inicializará los periféricos y entregará
        magnitudes en unidades acordadas: velocidad angular en rad/s,
        distancia en metros, ángulos de rueda y señales de mando para motores.
        Incluirá HAL para tiempo, MPU6050, VL53L0X, AS5600, GPIO/PWM del DRV8833
        y telemetría. No realizará integración de pose, mapeo ni navegación.
  \item[Capa 2 --- drivers compartidos.] Interpretará las señales físicas de
        cada dispositivo y entregará datos listos para la aplicación. Incluirá
        la calibración y corrección del giroscopio, validación de distancia,
        lectura y conversión de ángulos AS5600, y generación de señales de
        control para el DRV8833. Sus interfaces no dependerán de Arduino ni de
        SAPI.
  \item[Capa 3 --- aplicación.] Coordinará las tareas FreeRTOS y contendrá la
        lógica propia del rover: adquisición y sincronización de muestras,
        cinemática diferencial y estimación de pose, actualización de la grilla
        de ocupación, selección de consignas según el modo y distribución de
        datos a la navegación y las comunicaciones. Los algoritmos de mapa y
        navegación no se ubicarán en la HAL ni en los drivers de dispositivos.
\end{description}

El runtime común organizará el trabajo en tareas de adquisición, procesamiento
de drivers/aplicación y publicación de telemetría. La división interna de la
tarea de procesamiento podrá refinarse al medir los tiempos de ejecución; el
contrato entre capas se mantendrá igual para las dos placas.

\subsection{Bibliotecas y herramientas externas previstas}

\begin{itemize}
  \item \textbf{ESP32:} framework Arduino y su integración con FreeRTOS;
        \texttt{Wire} para I2C; biblioteca Adafruit \texttt{VL53L0X} para el
        sensor ToF. Para el transporte inalámbrico se utilizará la capacidad
        disponible en ESP32 una vez elegido el protocolo de enlace.
  \item \textbf{EDU-CIAA:} LPCOpen y SAPI para la placa y sus periféricos, más
        el kernel FreeRTOS configurado en el proyecto. La HAL de ToF implementa
        el acceso I2C por registros y no utiliza la biblioteca Adafruit.
  \item \textbf{Código del proyecto:} contratos HAL, drivers de MPU6050,
        VL53L0X, AS5600 y DRV8833, y módulos de pose, mapa y navegación. Los
        drivers de AS5600 y DRV8833 se incorporarán al contar con el hardware y
        sus interfaces acordadas.
  \item \textbf{PC:} Processing y su biblioteca Serial se utilizan en V0 para
        visualizar las muestras recibidas. La interfaz objetivo mostrará la
        pose del rover y el mapa actualizado.
\end{itemize}

\subsection{Software de PC e interfaz de usuario prevista}

El alcance confirmado para el programa de PC es visualizar la pose del rover y
el mapa actualizado. El enlace serie USB/UART se utiliza en V0; para transmitir
pose y mapa sin cable se prevé incorporar Wi-Fi o Bluetooth en una etapa
posterior. La elección entre ambos protocolos queda pendiente. Los
requerimientos generales también mencionan órdenes START/STOP y, como objetivos
secundarios, teleoperación y planificación de rutas; su incorporación como
controles de esta interfaz queda por definir.

\begin{center}
\fbox{\parbox{0.91\linewidth}{\centering
Inicio $\rightarrow$ conexión con el rover $\rightarrow$ espera de datos
$\rightarrow$ visualización continua de pose y mapa $\rightarrow$ pérdida de
enlace: indicar desconexión / conexión restablecida: reanudar visualización}}
\end{center}

Si se pierde el enlace durante la operación, el firmware deberá aplicar la
parada segura y la interfaz indicará la desconexión. El flujo representa la
interfaz prevista; las pantallas de Processing disponibles actualmente solo
reciben y dibujan telemetría.

\subsection{Ensayos realizados y estado de validación}

Los ensayos se realizaron con el firmware V0 y hardware conectado, usando
programas de prueba y Processing; no se dispone de una simulación integral del
rover. Se separaron perfiles de prueba para compilar la lectura de IMU y la del
ToF de manera independiente. También se implementó \texttt{imu\_compare}, que
transmite gyro Z crudo, ángulo crudo integrado y ángulo procesado para comparar
la integración sin corrección con la que aplica calibración y zona muerta.

El firmware V0 de ESP32 fue cargado y validado con Processing: se comprobó la
recepción de las muestras serie y la visualización de orientación y distancia.
El contrato de aplicación V0 transmite a 115200 baudios dos columnas,
\texttt{theta\_rad,distancia\_m}; una distancia inválida se representa con
$-1$. El build actual con FreeRTOS compartido compila, pero su telemetría debe
confirmarse después de la última modificación. El firmware EDU-CIAA compila y
enlaza, aunque falta probarlo con sensores conectados a la placa.

Los archivos disponibles no contienen capturas de pantalla ni resultados
numéricos de precisión o error. Para documentar visualmente los ensayos se debe
incorporar una captura de Processing durante la prueba V0 y, cuando se realice,
otra del gráfico de comparación de la IMU. Estas evidencias no se reemplazan
por imágenes ilustrativas.
.
% Este archivo contiene un fragmento del informe, no un documento independiente.

% =============================================================================
\section{Diseño del firmware, simulación y depuración}
\label{sec:diseno-firmware}
% =============================================================================

\subsection{Arquitectura propuesta y criterios de diseño}

El firmware se diseñará para integrar la percepción, la estimación de la pose,
la construcción del mapa y el control del rover diferencial. La arquitectura
será portable entre plataformas: el desarrollo y la validación funcional se
realizarán primero sobre ESP32; luego se portará a EDU-CIAA reutilizando la
aplicación y los drivers compartidos, y reemplazando la capa de adaptación al
hardware.

Se adopta una organización en tres capas para separar las operaciones propias
de cada placa de los algoritmos del robot. Esta separación permite probar los
módulos por partes, mantener una única implementación de la lógica del rover y
evitar que las decisiones de una placa condicionen el código común. La
arquitectura de ejecución se organizará en tareas periódicas controladas por
tiempo y planificadas por FreeRTOS. El scheduler será apropiativo y asignará el
procesador según las prioridades; las tareas se bloquearán durante sus
intervalos de espera. Por lo tanto, el modelo es de tareas periódicas con
planificación apropiativa, no un superloop cooperativo. Se utilizará FreeRTOS
en ambas plataformas.

\begin{center}
\fbox{\parbox{0.96\linewidth}{\centering
Sensores y actuadores $\leftrightarrow$ Capa 1: HAL de la placa
$\leftrightarrow$ Capa 2: drivers de sensores y motores
$\leftrightarrow$ Capa 3: adquisición, pose, mapa y navegación
$\leftrightarrow$ comunicaciones $\leftrightarrow$ interfaz de PC}}
\end{center}

El modelo elegido prioriza interfaces estables entre capas, unidades físicas
comunes y un runtime compartido. El período de muestreo y las prioridades se
ajustarán mediante mediciones de carga y latencia cuando se integren los
encoders, la construcción del mapa y el control de motores.

\subsection{Pseudocódigo del funcionamiento previsto}

\begin{enumerate}
  \item Configurar la placa, el scheduler, el canal de telemetría y el estado
        inicial seguro de los actuadores.
  \item Inicializar la IMU MPU6050, el sensor ToF VL53L0X, los encoders AS5600
        y el puente H DRV8833. Informar qué periféricos no están disponibles;
        no habilitar el movimiento si falta una condición necesaria para el
        control seguro.
  \item Calibrar los sensores que lo requieran y poner en cero la pose inicial
        $(x,y,\theta)$. Preparar la representación del mapa correspondiente:
        acumulación polar para V0 o grilla de ocupación para V1.
  \item Iniciar las tareas periódicas de FreeRTOS para adquisición,
        procesamiento de la aplicación y publicación/recepción de datos.
  \item En cada ciclo de adquisición, leer gyro Z, distancia ToF y ángulos de
        ambas ruedas; asociar las muestras con el instante de lectura y
        entregarlas a la aplicación.
  \item Procesar las muestras: corregir el gyro, estimar el desplazamiento de
        cada rueda, actualizar la pose del rover y validar la distancia.
  \item Actualizar el mapa con el rayo medido por el ToF y la pose vigente.
        En V0 se acumularán mediciones en un plot polar suponiendo que el rover
        gira en el origen; en V1 se actualizará una grilla de ocupación 2D en
        coordenadas mundo usando $(x,y,\theta)$.
  \item Según el modo seleccionado, obtener la consigna desde navegación
        autónoma o desde teleoperación, convertirla en velocidades de rueda y
        comandar el DRV8833 mediante la HAL. Ante una condición de parada o una
        pérdida de comunicación que supere el límite definido, detener los
        motores.
  \item Publicar pose, estado y datos del mapa para la interfaz externa;
        repetir el ciclo mientras el sistema esté activo.
\end{enumerate}

La navegación autónoma, el control de motores, la integración de AS5600 y la
grilla V1 forman parte del desarrollo previsto; no se presentan como
funcionalidades ya verificadas. SLAM permanece como objetivo secundario: la
primera grilla se construirá con la pose estimada por odometría.

\subsection{Módulos organizados en capas}

\begin{description}
  \item[Capa 1 --- HAL de periféricos.] Tendrá una interfaz común y una
        implementación por placa. Inicializará los periféricos y entregará
        magnitudes en unidades acordadas: velocidad angular en rad/s,
        distancia en metros, ángulos de rueda y señales de mando para motores.
        Incluirá HAL para tiempo, MPU6050, VL53L0X, AS5600, GPIO/PWM del DRV8833
        y telemetría. No realizará integración de pose, mapeo ni navegación.
  \item[Capa 2 --- drivers compartidos.] Interpretará las señales físicas de
        cada dispositivo y entregará datos listos para la aplicación. Incluirá
        la calibración y corrección del giroscopio, validación de distancia,
        lectura y conversión de ángulos AS5600, y generación de señales de
        control para el DRV8833. Sus interfaces no dependerán de Arduino ni de
        SAPI.
  \item[Capa 3 --- aplicación.] Coordinará las tareas FreeRTOS y contendrá la
        lógica propia del rover: adquisición y sincronización de muestras,
        cinemática diferencial y estimación de pose, actualización de la grilla
        de ocupación, selección de consignas según el modo y distribución de
        datos a la navegación y las comunicaciones. Los algoritmos de mapa y
        navegación no se ubicarán en la HAL ni en los drivers de dispositivos.
\end{description}

El runtime común organizará el trabajo en tareas de adquisición, procesamiento
de drivers/aplicación y publicación de telemetría. La división interna de la
tarea de procesamiento podrá refinarse al medir los tiempos de ejecución; el
contrato entre capas se mantendrá igual para las dos placas.

\subsection{Bibliotecas y herramientas externas previstas}

\begin{itemize}
  \item \textbf{ESP32:} framework Arduino y su integración con FreeRTOS;
        \texttt{Wire} para I2C; biblioteca Adafruit \texttt{VL53L0X} para el
        sensor ToF. Para el transporte inalámbrico se utilizará la capacidad
        disponible en ESP32 una vez elegido el protocolo de enlace.
  \item \textbf{EDU-CIAA:} LPCOpen y SAPI para la placa y sus periféricos, más
        el kernel FreeRTOS configurado en el proyecto. La HAL de ToF implementa
        el acceso I2C por registros y no utiliza la biblioteca Adafruit.
  \item \textbf{Código del proyecto:} contratos HAL, drivers de MPU6050,
        VL53L0X, AS5600 y DRV8833, y módulos de pose, mapa y navegación. Los
        drivers de AS5600 y DRV8833 se incorporarán al contar con el hardware y
        sus interfaces acordadas.
  \item \textbf{PC:} Processing y su biblioteca Serial se utilizan en V0 para
        visualizar las muestras recibidas. La interfaz objetivo mostrará la
        pose del rover y el mapa actualizado.
\end{itemize}

\subsection{Software de PC e interfaz de usuario prevista}

El alcance confirmado para el programa de PC es visualizar la pose del rover y
el mapa actualizado. El enlace serie USB/UART se utiliza en V0; para transmitir
pose y mapa sin cable se prevé incorporar Wi-Fi o Bluetooth en una etapa
posterior. La elección entre ambos protocolos queda pendiente. Los
requerimientos generales también mencionan órdenes START/STOP y, como objetivos
secundarios, teleoperación y planificación de rutas; su incorporación como
controles de esta interfaz queda por definir.

\begin{center}
\fbox{\parbox{0.91\linewidth}{\centering
Inicio $\rightarrow$ conexión con el rover $\rightarrow$ espera de datos
$\rightarrow$ visualización continua de pose y mapa $\rightarrow$ pérdida de
enlace: indicar desconexión / conexión restablecida: reanudar visualización}}
\end{center}

Si se pierde el enlace durante la operación, el firmware deberá aplicar la
parada segura y la interfaz indicará la desconexión. El flujo representa la
interfaz prevista; las pantallas de Processing disponibles actualmente solo
reciben y dibujan telemetría.

\subsection{Ensayos realizados y estado de validación}

Los ensayos se realizaron con el firmware V0 y hardware conectado, usando
programas de prueba y Processing; no se dispone de una simulación integral del
rover. Se separaron perfiles de prueba para compilar la lectura de IMU y la del
ToF de manera independiente. También se implementó \texttt{imu\_compare}, que
transmite gyro Z crudo, ángulo crudo integrado y ángulo procesado para comparar
la integración sin corrección con la que aplica calibración y zona muerta.

El firmware V0 de ESP32 fue cargado y validado con Processing: se comprobó la
recepción de las muestras serie y la visualización de orientación y distancia.
El contrato de aplicación V0 transmite a 115200 baudios dos columnas,
\texttt{theta\_rad,distancia\_m}; una distancia inválida se representa con
$-1$. El build actual con FreeRTOS compartido compila, pero su telemetría debe
confirmarse después de la última modificación. El firmware EDU-CIAA compila y
enlaza, aunque falta probarlo con sensores conectados a la placa.

Los archivos disponibles no contienen capturas de pantalla ni resultados
numéricos de precisión o error. Para documentar visualmente los ensayos se debe
incorporar una captura de Processing durante la prueba V0 y, cuando se realice,
otra del gráfico de comparación de la IMU. Estas evidencias no se reemplazan
por imágenes ilustrativas.
.
% Este archivo contiene un fragmento del informe, no un documento independiente.

% =============================================================================
\section{Diseño del firmware, simulación y depuración}
\label{sec:diseno-firmware}
% =============================================================================

\subsection{Arquitectura propuesta y criterios de diseño}

El firmware se diseñará para integrar la percepción, la estimación de la pose,
la construcción del mapa y el control del rover diferencial. La arquitectura
será portable entre plataformas: el desarrollo y la validación funcional se
realizarán primero sobre ESP32; luego se portará a EDU-CIAA reutilizando la
aplicación y los drivers compartidos, y reemplazando la capa de adaptación al
hardware.

Se adopta una organización en tres capas para separar las operaciones propias
de cada placa de los algoritmos del robot. Esta separación permite probar los
módulos por partes, mantener una única implementación de la lógica del rover y
evitar que las decisiones de una placa condicionen el código común. La
arquitectura de ejecución se organizará en tareas periódicas controladas por
tiempo y planificadas por FreeRTOS. El scheduler será apropiativo y asignará el
procesador según las prioridades; las tareas se bloquearán durante sus
intervalos de espera. Por lo tanto, el modelo es de tareas periódicas con
planificación apropiativa, no un superloop cooperativo. Se utilizará FreeRTOS
en ambas plataformas.

\begin{center}
\fbox{\parbox{0.96\linewidth}{\centering
Sensores y actuadores $\leftrightarrow$ Capa 1: HAL de la placa
$\leftrightarrow$ Capa 2: drivers de sensores y motores
$\leftrightarrow$ Capa 3: adquisición, pose, mapa y navegación
$\leftrightarrow$ comunicaciones $\leftrightarrow$ interfaz de PC}}
\end{center}

El modelo elegido prioriza interfaces estables entre capas, unidades físicas
comunes y un runtime compartido. El período de muestreo y las prioridades se
ajustarán mediante mediciones de carga y latencia cuando se integren los
encoders, la construcción del mapa y el control de motores.

\subsection{Pseudocódigo del funcionamiento previsto}

\begin{enumerate}
  \item Configurar la placa, el scheduler, el canal de telemetría y el estado
        inicial seguro de los actuadores.
  \item Inicializar la IMU MPU6050, el sensor ToF VL53L0X, los encoders AS5600
        y el puente H DRV8833. Informar qué periféricos no están disponibles;
        no habilitar el movimiento si falta una condición necesaria para el
        control seguro.
  \item Calibrar los sensores que lo requieran y poner en cero la pose inicial
        $(x,y,\theta)$. Preparar la representación del mapa correspondiente:
        acumulación polar para V0 o grilla de ocupación para V1.
  \item Iniciar las tareas periódicas de FreeRTOS para adquisición,
        procesamiento de la aplicación y publicación/recepción de datos.
  \item En cada ciclo de adquisición, leer gyro Z, distancia ToF y ángulos de
        ambas ruedas; asociar las muestras con el instante de lectura y
        entregarlas a la aplicación.
  \item Procesar las muestras: corregir el gyro, estimar el desplazamiento de
        cada rueda, actualizar la pose del rover y validar la distancia.
  \item Actualizar el mapa con el rayo medido por el ToF y la pose vigente.
        En V0 se acumularán mediciones en un plot polar suponiendo que el rover
        gira en el origen; en V1 se actualizará una grilla de ocupación 2D en
        coordenadas mundo usando $(x,y,\theta)$.
  \item Según el modo seleccionado, obtener la consigna desde navegación
        autónoma o desde teleoperación, convertirla en velocidades de rueda y
        comandar el DRV8833 mediante la HAL. Ante una condición de parada o una
        pérdida de comunicación que supere el límite definido, detener los
        motores.
  \item Publicar pose, estado y datos del mapa para la interfaz externa;
        repetir el ciclo mientras el sistema esté activo.
\end{enumerate}

La navegación autónoma, el control de motores, la integración de AS5600 y la
grilla V1 forman parte del desarrollo previsto; no se presentan como
funcionalidades ya verificadas. SLAM permanece como objetivo secundario: la
primera grilla se construirá con la pose estimada por odometría.

\subsection{Módulos organizados en capas}

\begin{description}
  \item[Capa 1 --- HAL de periféricos.] Tendrá una interfaz común y una
        implementación por placa. Inicializará los periféricos y entregará
        magnitudes en unidades acordadas: velocidad angular en rad/s,
        distancia en metros, ángulos de rueda y señales de mando para motores.
        Incluirá HAL para tiempo, MPU6050, VL53L0X, AS5600, GPIO/PWM del DRV8833
        y telemetría. No realizará integración de pose, mapeo ni navegación.
  \item[Capa 2 --- drivers compartidos.] Interpretará las señales físicas de
        cada dispositivo y entregará datos listos para la aplicación. Incluirá
        la calibración y corrección del giroscopio, validación de distancia,
        lectura y conversión de ángulos AS5600, y generación de señales de
        control para el DRV8833. Sus interfaces no dependerán de Arduino ni de
        SAPI.
  \item[Capa 3 --- aplicación.] Coordinará las tareas FreeRTOS y contendrá la
        lógica propia del rover: adquisición y sincronización de muestras,
        cinemática diferencial y estimación de pose, actualización de la grilla
        de ocupación, selección de consignas según el modo y distribución de
        datos a la navegación y las comunicaciones. Los algoritmos de mapa y
        navegación no se ubicarán en la HAL ni en los drivers de dispositivos.
\end{description}

El runtime común organizará el trabajo en tareas de adquisición, procesamiento
de drivers/aplicación y publicación de telemetría. La división interna de la
tarea de procesamiento podrá refinarse al medir los tiempos de ejecución; el
contrato entre capas se mantendrá igual para las dos placas.

\subsection{Bibliotecas y herramientas externas previstas}

\begin{itemize}
  \item \textbf{ESP32:} framework Arduino y su integración con FreeRTOS;
        \texttt{Wire} para I2C; biblioteca Adafruit \texttt{VL53L0X} para el
        sensor ToF. Para el transporte inalámbrico se utilizará la capacidad
        disponible en ESP32 una vez elegido el protocolo de enlace.
  \item \textbf{EDU-CIAA:} LPCOpen y SAPI para la placa y sus periféricos, más
        el kernel FreeRTOS configurado en el proyecto. La HAL de ToF implementa
        el acceso I2C por registros y no utiliza la biblioteca Adafruit.
  \item \textbf{Código del proyecto:} contratos HAL, drivers de MPU6050,
        VL53L0X, AS5600 y DRV8833, y módulos de pose, mapa y navegación. Los
        drivers de AS5600 y DRV8833 se incorporarán al contar con el hardware y
        sus interfaces acordadas.
  \item \textbf{PC:} Processing y su biblioteca Serial se utilizan en V0 para
        visualizar las muestras recibidas. La interfaz objetivo mostrará la
        pose del rover y el mapa actualizado.
\end{itemize}

\subsection{Software de PC e interfaz de usuario prevista}

El alcance confirmado para el programa de PC es visualizar la pose del rover y
el mapa actualizado. El enlace serie USB/UART se utiliza en V0; para transmitir
pose y mapa sin cable se prevé incorporar Wi-Fi o Bluetooth en una etapa
posterior. La elección entre ambos protocolos queda pendiente. Los
requerimientos generales también mencionan órdenes START/STOP y, como objetivos
secundarios, teleoperación y planificación de rutas; su incorporación como
controles de esta interfaz queda por definir.

\begin{center}
\fbox{\parbox{0.91\linewidth}{\centering
Inicio $\rightarrow$ conexión con el rover $\rightarrow$ espera de datos
$\rightarrow$ visualización continua de pose y mapa $\rightarrow$ pérdida de
enlace: indicar desconexión / conexión restablecida: reanudar visualización}}
\end{center}

Si se pierde el enlace durante la operación, el firmware deberá aplicar la
parada segura y la interfaz indicará la desconexión. El flujo representa la
interfaz prevista; las pantallas de Processing disponibles actualmente solo
reciben y dibujan telemetría.

\subsection{Ensayos realizados y estado de validación}

Los ensayos se realizaron con el firmware V0 y hardware conectado, usando
programas de prueba y Processing; no se dispone de una simulación integral del
rover. Se separaron perfiles de prueba para compilar la lectura de IMU y la del
ToF de manera independiente. También se implementó \texttt{imu\_compare}, que
transmite gyro Z crudo, ángulo crudo integrado y ángulo procesado para comparar
la integración sin corrección con la que aplica calibración y zona muerta.

El firmware V0 de ESP32 fue cargado y validado con Processing: se comprobó la
recepción de las muestras serie y la visualización de orientación y distancia.
El contrato de aplicación V0 transmite a 115200 baudios dos columnas,
\texttt{theta\_rad,distancia\_m}; una distancia inválida se representa con
$-1$. El build actual con FreeRTOS compartido compila, pero su telemetría debe
confirmarse después de la última modificación. El firmware EDU-CIAA compila y
enlaza, aunque falta probarlo con sensores conectados a la placa.

Los archivos disponibles no contienen capturas de pantalla ni resultados
numéricos de precisión o error. Para documentar visualmente los ensayos se debe
incorporar una captura de Processing durante la prueba V0 y, cuando se realice,
otra del gráfico de comparación de la IMU. Estas evidencias no se reemplazan
por imágenes ilustrativas.
